\section{引言：什么是计算机视觉}

视觉是人类最重要的感官之一。它使我们能够检测和解读情感与面部表情、在世界中导航、操纵物体以及与周围环境互动。Alexander Amini 在这节课的开篇提出了一个简洁的定义：计算机视觉（Computer Vision）就是赋予计算机"通过观察来了解什么在哪里"（\textit{to know what is where by looking}）的能力。

\begin{figure}[H]
\centering
\includegraphics[width=0.85\textwidth]{slides-images/slide-02.jpg}
\caption{计算机视觉的核心定义："To know what is where by looking." \protect\footnotemark}
\end{figure}
\footnotetext{Slides 第2页。}

视觉不仅仅是静态地识别物体，还包括理解场景的动态变化。例如，当我们看到一张城市街道照片时，我们不仅能识别出汽车、行人和交通信号灯，还能推断出哪些汽车在行驶、哪些停在路边，以及行人的运动方向。这种从单幅图像中推断动态信息的能力，正是计算机视觉希望实现的目标。

\begin{figure}[H]
\centering
\includegraphics[width=0.85\textwidth]{slides-images/slide-03.jpg}
\caption{计算机视觉的目标：从图像中发现什么物体在场景中、它们在哪里、正在发生什么动作，以及预测和推断未来事件。\protect\footnotemark}
\end{figure}
\footnotetext{Slides 第3页。}

\begin{knowledgebox}{为什么计算机视觉如此困难？}
人类视觉系统经过数百万年的进化，能够毫不费力地完成物体识别、场景理解等任务。但对计算机而言，即使是同一物体，由于视角变化（viewpoint variation）、尺度变化（scale variation）、形变（deformation）、遮挡（occlusion）、光照变化（illumination conditions）、背景杂乱（background clutter）和类内变化（intra-class variation）等因素，在像素层面的表现可能截然不同。这些都使得计算机视觉成为一个极具挑战性的问题。
\end{knowledgebox}

\subsection{计算机视觉的应用与影响}

深度学习正在推动计算机视觉领域的巨大革命，其应用无处不在：

\begin{figure}[H]
\centering
\includegraphics[width=0.9\textwidth]{slides-images/slide-04.jpg}
\caption{计算机视觉的广泛应用：机器人、无障碍辅助、生物医学、自动驾驶和移动计算。\protect\footnotemark}
\end{figure}
\footnotetext{Slides 第4页。}

\begin{itemize}
    \item \textbf{面部检测与识别}：不仅检测人脸的存在，还能分析面部关键点和微表情。这是最早获得广泛应用的计算机视觉任务之一。
    \item \textbf{自动驾驶}：MIT 的 Amini 实验室曾构建端到端自动驾驶系统，仅凭视觉输入即可在从未见过的道路上自主驾驶。与大多数自动驾驶汽车使用模块化流水线（多个模型加预定义地图）不同，端到端方法使用单一模型直接从图像到方向盘转角。
    \item \textbf{医学影像}：CNN 在乳腺癌筛查中已超越专业放射科医生的准确率，能够检测到人类医生遗漏的病灶。此外还广泛应用于皮肤癌检测、COVID-19 影像诊断等领域。
    \item \textbf{无障碍辅助}：例如 Google Project Guideline 项目，帮助视障人士独立跑步。
    \item \textbf{移动计算}：每个人的手机中都运行着大量计算机视觉模型。
\end{itemize}

\begin{figure}[H]
\centering
\begin{subfigure}[b]{0.48\textwidth}
\centering
\includegraphics[width=\textwidth]{slides-images/slide-05.jpg}
\caption{面部检测与关键点识别}
\end{subfigure}
\hfill
\begin{subfigure}[b]{0.48\textwidth}
\centering
\includegraphics[width=\textwidth]{slides-images/slide-07.jpg}
\caption{医学影像中的计算机视觉}
\end{subfigure}
\caption{计算机视觉的代表性应用。\protect\footnotemark}
\end{figure}
\footnotetext{Slides 第5页和第7页。}

\subsection{本章小结}

计算机视觉是赋予计算机理解视觉世界的能力。虽然人类做起来轻而易举，但由于物体外观的巨大变异性，这对计算机来说极具挑战。深度学习的兴起正在从根本上改变这一领域，使得计算机在许多视觉任务上达到甚至超越人类水平。

\newpage
